This exercise focuses on the binomial one-period model, which is the foundational building block for pricing options. We have two states of the world at $t=1$: a ``down'' state ($\omega_1$) and an ``up'' state ($\omega_2$).

\subsection*{Question 1: Arbitrage-free condition}
We need to prove that the market is arbitrage-free if and only if the risk-free gross return $1+r$ lies strictly between the down-factor $d$ and the up-factor $u$. We can prove this by examining the three possible scenarios for $1+r$.

\textbf{Case 1: Suppose $1+r \le d$} \\
In this scenario, the risk-free rate is so low that even the worst-case return of the stock ($d$) is better than or equal to the bank's interest rate. 
We can construct an arbitrage strategy by borrowing money from the bank to buy the stock. 
Let our strategy be $(\xi^0, \xi) = (-S_0, 1)$. This means we borrow $S_0$ from the bank (so our cash position is $-S_0$) and buy 1 share of the stock.
\begin{itemize}
    \item Initial cost: $V_0 = -S_0 + 1 \cdot S_0 = 0$.
    \item Value at $t=1$ in the worst case (down state): 
    \begin{equation}
        V_1(\omega_1) = -S_0(1+r) + S_0 d = S_0(d - (1+r)) \ge 0
    \end{equation}
    \item Value at $t=1$ in the best case (up state):
    \begin{equation}
        V_1(\omega_2) = -S_0(1+r) + S_0 u = S_0(u - (1+r)) > 0
    \end{equation}
\end{itemize}
Since $V_0 = 0$, $V_1 \ge 0$ always, and $V_1 > 0$ with strictly positive probability, this is an arbitrage. Therefore, for the market to be arbitrage-free, $1+r$ must be strictly greater than $d$.

\textbf{Case 2: Suppose $1+r \ge u$} \\
Here, the bank offers a guaranteed return that is better than or equal to the best-case scenario of the stock.
We can construct an arbitrage by short-selling the stock and putting the money in the bank. Let our strategy be $(\xi^0, \xi) = (S_0, -1)$.
\begin{itemize}
    \item Initial cost: $V_0 = S_0 - 1 \cdot S_0 = 0$.
    \item Value at $t=1$ in the best case (up state):
    \begin{equation}
        V_1(\omega_2) = S_0(1+r) - S_0 u = S_0((1+r) - u) \ge 0
    \end{equation}
    \item Value at $t=1$ in the worst case (down state):
    \begin{equation}
        V_1(\omega_1) = S_0(1+r) - S_0 d = S_0((1+r) - d) > 0
    \end{equation}
\end{itemize}
Again, this guarantees profits \& avoids losses without any initial investment, which is an arbitrage. Thus, we must have $1+r < u$.

\textbf{Case 3: Suppose $1+r \in (d,u)$} \\
By the First Fundamental Theorem of Asset Pricing, a market is arbitrage-free if and only if there exists an equivalent martingale measure (EMM) denoted by $Q$. Under $Q$, the discounted stock price must be a martingale:
\begin{equation}
    E^Q\left[\frac{S_1}{1+r}\right] = S_0
\end{equation}
Let $q = Q(\{\omega_2\})$ be the risk-neutral probability of the up state. Then $1-q = Q(\{\omega_1\})$ is the probability of the down state. Expanding the expectation:
\begin{equation}
    \frac{1}{1+r} \left( (1-q)S_0 d + q S_0 u \right) = S_0
\end{equation}
Divide both sides by $S_0$ and multiply by $1+r$:
\begin{equation}
    d - qd + qu = 1+r \implies q(u-d) = 1+r-d \implies q = \frac{1+r-d}{u-d}
\end{equation}
For $Q$ to be a valid equivalent probability measure, we must have $0 < q < 1$. 
Since $1+r \in (d,u)$, we know $1+r-d > 0$ and $u-d > 0$, so $q > 0$. We also know $1+r < u$, meaning $1+r-d < u-d$, which gives $q < 1$. Thus, a valid EMM exists, proving the market is arbitrage-free.

\subsection*{Question 2: Completeness of the market}
The Second Fundamental Theorem of Asset Pricing states that an arbitrage-free market is complete if and only if the equivalent martingale measure $Q$ is \textit{unique}.

From Question 1, we established that avoiding arbitrage requires $1+r \in (d,u)$. We then derived the risk-neutral probability for the up state:
\begin{equation}
    q = Q(\{\omega_2\}) = \frac{1+r-d}{u-d}
\end{equation}
Because $r$, $d$, and $u$ are fixed parameters of the market, there is exactly one mathematical solution for $q$. Since $Q$ is unique, every derivative can be perfectly replicated by a portfolio of the underlying stock and the riskless bond. Therefore, the market is complete.

\subsection*{Question 3: Call option price}
In a complete, arbitrage-free market, the unique price of any derivative $\pi_0$ is the discounted expected payoff under the risk-neutral measure $Q$:
\begin{equation}
    \pi_0 = E^Q\left[\frac{(S_1 - K)_+}{1+r}\right]
\end{equation}
Where $(x)_+ = \max(x, 0)$. We substitute our unique probabilities $q$ and $1-q$ into the expectation:
\begin{equation}
    \pi_0 = \frac{1}{1+r} \left( q \cdot (S_1(\omega_2) - K)_+ + (1-q) \cdot (S_1(\omega_1) - K)_+ \right)
\end{equation}
Substitute $q = \frac{1+r-d}{u-d}$ and $1-q = \frac{u-1-r}{u-d}$:
\begin{equation}
    \pi_0 = \frac{1+r-d}{u-d} \frac{(S_0 u - K)_+}{1+r} + \frac{u-1-r}{u-d} \frac{(S_0 d - K)_+}{1+r}
\end{equation}
This is the general pricing formula for a call option in a one-period binomial model.

\subsection*{Question 4: Numerical application}
We are given: $r = 0$, $S_0 = 100$, $u = 1.2$, $d = 0.9$, $K = 100$.

\textbf{1. Stock Returns ($\rho(S)$):}
The return is calculated as $\frac{S_1 - S_0}{S_0}$.
\begin{itemize}
    \item Up state: $S_1 = 100 \times 1.2 = 120$. Return $= \frac{120 - 100}{100} = 20\%$.
    \item Down state: $S_1 = 100 \times 0.9 = 90$. Return $= \frac{90 - 100}{100} = -10\%$.
\end{itemize}
So, $\rho(S) \in \{-10\%, 20\%\}$.

\textbf{2. Call Option Price ($\pi_0^C$):}
First, find the risk-neutral probability $q$:
\begin{equation}
    q = \frac{1+0-0.9}{1.2-0.9} = \frac{0.1}{0.3} = \frac{1}{3}
\end{equation}
Now, calculate the option payoffs:
\begin{itemize}
    \item Up state payoff: $(120 - 100)_+ = 20$
    \item Down state payoff: $(90 - 100)_+ = 0$
\end{itemize}
The price is the expected discounted payoff ($r=0$, so discount factor is 1):
\begin{equation}
    \pi_0^C = \frac{1}{3}(20) + \frac{2}{3}(0) = \frac{20}{3} \approx 6.67
\end{equation}

\textbf{3. Call Option Returns ($\rho(C)$):}
The return of the option is $\frac{\text{Payoff} - \text{Price}}{\text{Price}}$.
\begin{itemize}
    \item Up state return: $\frac{20 - 20/3}{20/3} = \frac{60/3 - 20/3}{20/3} = \frac{40/3}{20/3} = 2 = 200\%$.
    \item Down state return: $\frac{0 - 20/3}{20/3} = -1 = -100\%$.
\end{itemize}
So, $\rho(C) \in \{-100\%, 200\%\}$.

\textbf{Conclusion:} The initial investment required for the option (20/3) is much cheaper than the stock (100). However, the option has a massive leverage effect: it offers the potential for much higher percentage gains (200\% vs 20\%) but also total loss of the invested capital (-100\% vs -10\%).